\section{Introduction}

A visible graph records which local transitions are available. It need not
record what is retained when a path closes. A lifted system may return to the
same visible vertex while its receipt coordinate has changed. On a single
cycle, that retained receipt is one group element. On a general graph, there
are several independent circuits, and their receipts form a representation of
the graph fundamental group.

This paper develops that graph-wide algebra. The ingredients are deliberately
minimal: a connected finite graph $X$, a finite receipt group $R$, and one
inverse-compatible group element on every oriented edge. These data construct
a regular lift with vertex set $V(X)\times R$. The left action of $R$ changes
the receipt coordinate, while edge transport multiplies it on the right.
These actions commute even when $R$ is nonabelian.

Gauge makes the local edge table highly redundant. After choosing a spanning
tree, every tree receipt can be removed. The remaining receipts are exactly
the values on a free generating set of $\pi_1(X)$. Residual gauge conjugates
all generator images simultaneously. This gives the complete classification.

The image subgroup of the holonomy representation has a second role: it
determines reachability in the lift. Its cosets index connected components.
Thus classification and connectivity arise from the same circuit-receipt map.
The result belongs to the standard theory of voltage graphs and regular graph
covers \cite{gross-tucker}. Our purpose is to isolate it as a compact receipt
algebra and to state exactly what additional certificate is needed before it
can be applied to a raw finite graph.
